\section{Performance Evaluation and Metrics}
\label{sec:performance_evaluation}

To validate the viability of the DBaaS platform, performance evaluations were conducted on the local \texttt{dbaas-local} k3d cluster used throughout this thesis. Measurements cover provisioning latency, PGProxy routing behavior, and idle tenant memory footprint.

\subsection{Evaluation Methodology}

All benchmarks were executed against the deployment produced by \texttt{start.sh} on a single-node k3d cluster (K3s) with CloudNativePG and the MongoDB Community Operator pre-installed. Hardware and software context for the runs documented here:

\begin{itemize}
    \item \textbf{Cluster:} k3d cluster \texttt{dbaas-local}, one server node.
    \item \textbf{Tier under test:} \texttt{free} (0.5 CPU, 512 MiB memory request, 5 GiB storage).
    \item \textbf{Provisioning samples:} five consecutive PostgreSQL project creations; timer starts at successful \texttt{POST /api/v1/projects} response and ends when project \texttt{status} becomes \texttt{running} (polled every two seconds).
    \item \textbf{Query latency:} \texttt{pgbench} TPC-B-style workload (\texttt{-c 1 -j 1 -T 10}) executed inside the CNPG pod against database \texttt{app}.
    \item \textbf{Memory footprint:} \texttt{kubectl top pods} RSS after the tenant pods reached steady \texttt{Running} state.
\end{itemize}

Managed cloud databases such as AWS RDS or Google Cloud SQL typically require several minutes to allocate compute, configure networking, and bootstrap storage \cite{aws_rds_provisioning}.\footnote{Provisioning duration varies by instance class, region, and high-availability options; the comparison here is qualitative.} The measurements below characterize the platform's local operator-driven path rather than a formal cloud SLA benchmark.

\subsection{Database Provisioning Latency}

Five free-tier PostgreSQL projects were created sequentially through the REST API. Time-to-\texttt{running} included namespace creation, Secret injection, CloudNativePG \texttt{Cluster} reconciliation, PVC binding, pod startup, and PostgREST sidecar deployment.

\begin{table}[H]
    \centering
    \renewcommand{\arraystretch}{1.3}
    \begin{tabular}{|c|c|c|c|c|c|c|}
    \hline
    \textbf{Run} & 1 & 2 & 3 & 4 & 5 & \textbf{Mean} \\
    \hline
    Seconds to \texttt{running} & 20.4 & 26.5 & 26.4 & 30.5 & 30.5 & 26.9 \\
    \hline
    \end{tabular}
    \caption{Provisioning latency on local k3d cluster (free tier, PostgreSQL)}
    \label{tab:provisioning_latency}
\end{table}

The observed range was \textbf{20 to 31 seconds} with a mean of \textbf{26.9 seconds}. This includes secret generation, PVC attachment, database bootstrap, and PostgREST resource creation. Variation across runs reflects Kubernetes scheduling and storage binding jitter on a single-node development cluster.

\subsection{PGProxy Routing Overhead}

PGProxy accepts external PostgreSQL TCP connections, performs TLS \texttt{ClientHello} inspection to extract SNI, and forwards the byte stream without terminating TLS \cite{go_net_performance}. This preserves end-to-end encryption between the client and CloudNativePG, which is required for SCRAM channel-binding compatibility.

Inside the CNPG pod, a \texttt{pgbench} run with one client and a ten-second duration reported:

\begin{itemize}
    \item \textbf{Average transaction latency:} 0.834 ms
    \item \textbf{Throughput:} approximately 1{,}199 transactions per second
    \item \textbf{Initial connection time:} 48.8 ms (one-time setup cost)
\end{itemize}

An external \texttt{pgbench} comparison through PGProxy requires TPC-B tables owned by the tenant \texttt{app} role; the benchmark tables initialized by the superuser inside the pod are not accessible to external application credentials. For that reason, proxy overhead is characterized architecturally: after SNI resolution, PGProxy pipes bytes between client and database with no cryptographic termination step. Connection-level tests through the external hostname (\texttt{db-<id>.postgres.<domain>}) confirmed successful TLS passthrough routing. The dominant cost for short queries remains PostgreSQL execution and network RTT rather than proxy inspection.

\subsection{Resource Footprint of Idle Tenants}
\label{sec:resource_footprint}

Multi-tenant density depends on how much memory an idle project consumes. For one idle free-tier PostgreSQL project measured with \texttt{kubectl top pods}:

\begin{itemize}
    \item \textbf{CloudNativePG pod:} 86 MiB RSS
    \item \textbf{PostgREST sidecar:} 25 MiB RSS
    \item \textbf{Combined:} approximately 111 MiB per idle PostgreSQL project
\end{itemize}

On a worker node with 16 GiB of RAM, reserving roughly 2 GiB for the operating system, Kubernetes control-plane components, and database operators leaves about 14 GiB schedulable. At 111 MiB per idle stack, the node can host on the order of \textbf{120 to 130 idle free-tier PostgreSQL projects} before memory becomes the binding constraint. Actual density will be lower when projects execute queries, when storage caches grow, or when MongoDB projects (which run separate operator-managed replica-set pods) are mixed into the workload.

CPU limits configured in the resource tier table (0.5 cores for \texttt{free}) prevent a single tenant from monopolizing compute even when many idle projects are bin-packed onto the same node.
